% ============================================================================
% Chapter 37: Web Browsing and Searching
% ============================================================================
\chapter{Web Browsing and Searching}
\label{ch:browsing}

\begin{objectives}
  \item Use a web browser to navigate websites
  \item Search effectively using Google
  \item Bookmark favourite websites
  \item Evaluate website reliability and trustworthiness
\end{objectives}

\bytetip[tip]{Searching the web is a skill --- not all results are reliable, and knowing the right keywords makes you \textbf{10× faster} at finding what you need!}

\section{Using a Web Browser}
\label{sec:browser}
\index{web browser}

A \hl{web browser} is an application that lets you view websites. Popular browsers include \textbf{Google Chrome}, \textbf{Microsoft Edge}, \textbf{Mozilla Firefox}, and \textbf{Safari}.

\section{Searching Effectively}
\label{sec:searching}
\index{search engine}

\begin{theorybox}[Smart Search Tips]
\begin{enumerate}[label=\textcolor{ModuleBlue}{\textbf{\arabic*.}}, itemsep=3pt]
  \item \textbf{Use keywords}, not full sentences --- ``Indian Independence Day date'' instead of ``When is Indian Independence Day?''
  \item \textbf{Use quotes} for exact phrases --- \texttt{"solar system planets"}
  \item \textbf{Add \texttt{site:}} to search within one site --- \texttt{site:wikipedia.org dinosaurs}
  \item \textbf{Use minus} to exclude words --- \texttt{jaguar -car} (to find the animal)
  \item \textbf{Check multiple sources} --- don't trust just one website
\end{enumerate}
\end{theorybox}

\section{Evaluating Websites}
\label{sec:evaluate}
\index{website!reliability}

Not everything on the internet is true! Use the \textbf{CRAAP test} to check reliability:

\begin{table}[H]
\centering
\caption{The CRAAP test for evaluating websites}
\label{tab:craap}
\renewcommand{\arraystretch}{1.3}
\begin{tabularx}{\textwidth}{l l X}
\toprule
\rowcolor{HeaderRow}
\textcolor{white}{\textbf{Letter}} & \textcolor{white}{\textbf{Stands For}} & \textcolor{white}{\textbf{Ask Yourself}} \\
\midrule
\rowcolor{RowLight} C & Currency & Is the information up to date? \\
\rowcolor{RowDark} R & Relevance & Does it answer your question? \\
\rowcolor{RowLight} A & Authority & Who wrote it? Are they an expert? \\
\rowcolor{RowDark} A & Accuracy & Can you verify the facts elsewhere? \\
\rowcolor{RowLight} P & Purpose & Why was it written? To inform, sell, or persuade? \\
\bottomrule
\end{tabularx}
\end{table}

\bytetip[warn]{Anyone can publish anything on the internet. Just because it looks professional doesn't mean it's true. Always verify important facts with at least \textbf{two reliable sources}!}

\begin{funfact}
Google processes over \textbf{8.5 billion searches per day}. That's about 99,000 searches every single second! The most Googled question ever is ``What to watch?''
\end{funfact}

\begin{reviewqs}
  \item Name three popular web browsers.
  \item What is a search engine and name two examples.
  \item List three tips for searching effectively on Google.
  \item What is the CRAAP test?
  \item Why should you check multiple sources for information?
\end{reviewqs}

\begin{challenge}
Search for the answer to this question: ``How tall is the Eiffel Tower in metres?'' Use \textbf{three different websites} to find the answer. Do they all agree? Write down each source and its answer. Which source do you trust most, and why?
\end{challenge}
